\section{Metatheory}
\label{sec:metatheory}

This section summarizes the main metatheoretical results for \thecalculus. All
definitions from the previous sections have been formalized in
Agda. We mark the theorems we proved mechanically with a little agda bird after
the theorem name. Non-mechanized proofs are shown in the supplement.

We first state type safety for declarative typing, then weak operational
correspondence for the translation, and finally soundness and completeness of
algorithmic typing.

\subsection{Type Safety for Declarative Typing}

Preservation
yields the usual safety invariant across reduction steps.

\begin{theorem}[Expression preservation\;\AgdaBird]
  If $\HasType \SessCtx \E \T$ and $\E \EReduce \E*$, then
  $\HasType \SessCtx {\E*} \T$.
\end{theorem}

\begin{theorem}[Process preservation]
  If $\ProcHasType \SessCtx \Proc$ and $\Proc \PReduce \Proc*$, then
  $\ProcHasType \SessCtx {\Proc*}$.
\end{theorem}

We complement preservation with progress.
For expressions, progress includes a concurrency-specific blocked case. Informally,
an expression is blocked when evaluation reaches a communication primitive in an
evaluation context that cannot proceed on its own because it requires a matching
action from another thread. The formal blocked-expression predicate is given in \cref{fig:blocked-expr-proc}. We write $\SessCtx$
for a context $\Ctx$ that only contains session types.

\begin{theorem}[Expression progress\;\AgdaBird]
  \label{thm:expression-progress}
  If $\HasType \SessCtx \E \T$, then either $\E$ is a value or $\E$ is blocked on concurrent actions or there exists an $\E*$ such that $\E \EReduce \E*$.
\end{theorem}

A process is concurrently blocked, if the active threads are blocked and no
channel pair can synchronize. The formal blocked-process
predicate is given in \cref{fig:blocked-expr-proc}.
Process progress guarantees reduction only for non-blocked
configurations. Thus, deadlocks remain possible.

\begin{theorem}[Process progress]
  \label{thm:process-progress}
  If $\ProcHasType \SessCtx \Proc$, then either $\Proc \Congruent \PExp{\CUnit}$ or an expression in $\Proc$ is blocked on concurrent actions or
    there exists an $\Proc*$ such that $\Proc \PReduce \Proc*$.
\end{theorem}
The proof is deferred to \cref{sec:appendix-metatheory-ts}.

\subsection{Weak Simulation for Translation}

We relate source and target reduction through the process translation
$\TransP\Proc$, which maps a well-typed process $\ProcHasType \Ctx \Proc$ to its
run-time configuration under a value substitution $\Sub$ over the channel context
$\SessCtx$ of $\Ctx$. The translation introduces administrative
synchronization steps, hence both directions only hold up to a weak bisimilarity $\Bisim$,
which identifies run-time configurations that differ only by administrative
drop and acquire steps, discarded synchronization cells, and the reordering of
independent binders.

\begin{theorem}[Forward simulation\;\AgdaBird]
  If $\Proc \PReduce \Proc*$, then there exists a run-time configuration
  $\TProcQ$ such that $\TransP\Proc \TPReduce^{*} \TProcQ$ and
  $\TProcQ \Bisim \TransP{\Proc*}$.
\end{theorem}

The converse direction matches every target step by a finite source computation,
again up to $\Bisim$.

\begin{theorem}[Backward simulation]
  For all run-time configurations $\TProcR$ and $\TProcQ$ with
  $\TProcR \Bisim \TransP\Proc$, if $\TProcR \TPReduce \TProcQ$, then there
  exists a source process $\Proc*$ such that $\Proc \PReduce^{*} \Proc*$ and
  $\TProcQ \Bisim \TransP{\Proc*}$.
\end{theorem}
\subsection{Soundness and Completeness for Algorithmic Typing}

We now connect algorithmic typing with declarative typing. Soundness says
that any solved algorithmic judgment denotes a valid declarative typing, while
completeness says that declarative typings can be reconstructed algorithmically.

\begin{theorem}[Expression checking soundness\;\AgdaBird]
  If $\CheckType \E \AT \Eff$ and $\sigma \models \CSet$ is a closing
  substitution such that $\Fv(\AT) \subseteq \dom(\sigma)$, then
  $\HasType \Ctx \E {\AT{[\sigma]}}$.
\end{theorem}

The same soundness statement holds for inference mode.

\begin{theorem}[Expression inference soundness\;\AgdaBird]
  If $\InferType \E \AT \Eff$ and $\sigma \models \CSet$ is a closing
  substitution such that $\Fv(\AT) \subseteq \dom(\sigma)$, then
  $\HasType \Ctx \E {\AT{[\sigma]}}$.
\end{theorem}

For completeness, we again first cover checking and then synthesis.

\begin{theorem}[Expression checking completeness]
  If $\HasType \Ctx \E \T$, then
  $\CheckType \E \AT \Eff$ for some type $\AT$, set of constraints $\CSet$,
  and closing substitution $\sigma$ such that $\T = \AT{[\sigma]}$.
\end{theorem}

\begin{theorem}[Expression inference completeness]
  If $\HasType \Ctx \E \T$ and $\E$ is in synthesis form, then
  $\InferType \E \AT \Eff$ for some type $\AT$, set of constraints $\CSet$,
  and closing substitution $\sigma$ such that $\T = \AT{[\sigma]}$.
\end{theorem}

%%% Local Variables:
%%% mode: latex
%%% TeX-master: "../popl2027.tex"
%%% End:
